\scanpage{109}
% Archive foliation: 110. Reused unrelated typescript; printed page 18, n°382.
\begin{center}-- 18 --\hfill n\textsuperscript{o} 382\end{center}
\[
f(a+h)=\sum\Delta^s f(a).h^s\qquad\text{si $h$ est assez petit.}
\]
Soit $C_E^\omega(U)$ l'espace vectoriel des applications analytiques de $U$ dans $E$. Les $\Delta^s$ sont des endomorphismes $K$-linéaires de $C_E^\omega(U)$. On a
\[
\Delta^s\circ\Delta^t=\Delta^{s+t}.
\]
De plus, avec les notations et hypothèses du n\textsuperscript{o} 3.4.2., on a
\[
\Delta^s(f^1.f^2)=\sum_{u+v=s}\Delta^u f^1.\Delta^v f^2.
\]

3.5.3. Soit $f$ une application analytique d'un ouvert $U$ de $K^n$ dans $E$, et soit $g$ une application analytique d'un ouvert $V$ de $K^m$ dans $U$. L'application composée $f\circ g$ est analytique. En outre, supposons que $0\in V$ et que $g(0)=0$. Le développement de $f\circ g$ en série entière au point $0$ s'obtient alors en substituant, dans le développement de $f$ au point $0$, les développements en série entière des composantes $g_i$ de $g$ au point $0$ (cf. n\textsuperscript{o} 3.4.4.).

3.5.4. Soit $f=\sum b_s x^s$ une série convergente à coefficients dans $E$, et soit $D(f)$ son domaine de convergence stricte. La fonction définie par $f$ sur $D(f)$ est analytique. Pour tout $s\in\underline{\underline{N}}^n$, on a
\[
\Delta^s f(x)=\sum\binom{s+t}{s}b_{s+t}x^t,
\]
la série du membre de droite ayant un domaine de convergence stricte qui contient celui de $f$ (et lui est égal si $K=\underline{\underline{R}}$ ou $K=\underline{\underline{C}}$).

3.5.5. Si $f$ est analytique dans un ouvert $U$ de $K^n$, $f$ est \uline{strictement dérivable} dans $U$ (cf. n\textsuperscript{o} 1.2.4) et ses dérivées partielles $D_i f$ sont analytiques. Les dérivées partielles itérées $D^s f$ de $f$ sont données par la formule
\[
D^s f=s!\,\Delta^s f.
\]

3.5.6. Soient $f$ et $g$ deux applications analytiques d'un ouvert $U$ de $K^n$ dans $E$. Soit $a\in U$. Pour que $f$ et $g$ aient un contact d'ordre $\geq k$ en $a$, il faut et il suffit que $\Delta^s f(a)=\Delta^s g(a)$ si $|s|\leq k$. Pour que $f$ et $g$ aient un contact d'ordre infini en $a$, il faut et il suffit qu'elles soient égales dans un voisinage
